% =====================================================================
\chapter{Background}
\label{ch:background}
% =====================================================================

\section{Timing in Embedded Systems}
\label{sec:bg:timing}

Embedded firmware interacts with the physical world through deadlines:
a PWM edge must occur every control period, a sensor must be sampled at
a fixed rate, a bus frame must be acknowledged within a protocol
timeout. Two properties characterise the timing behaviour of such a
system:

\begin{description}
  \item[Accuracy] --- how close the \emph{mean} of an observed interval
        is to its nominal value (e.g.\ a \SI{10}{\milli\second} timer
        firing every \SI{10.001}{\milli\second} on average).
  \item[Jitter] --- the \emph{variance} of the interval around its
        mean, typically reported as a standard deviation or a
        percentile bound.
\end{description}

Real microcontrollers derive their notion of time from a crystal or RC
oscillator whose frequency exhibits static offset (specified in parts
per million, ppm), temperature drift, and cycle-to-cycle jitter. On top
of this physical layer, the RTOS adds software-induced timing effects:
tick quantisation, scheduling latency, priority inversion, and
interrupt masking windows. Any faithful emulation of firmware timing
must account for both layers.

\section{The Target Platform: STM32F103RB and Zephyr RTOS}
\label{sec:bg:platform}

All experiments in this thesis use the ST NUCLEO-F103RB development
board. Its STM32F103RB microcontroller~\cite{stm32rm0008} integrates an
Arm Cortex-M3 core~\cite{armcm3trm} clocked at \SI{72}{\mega\hertz},
\SI{20}{\kilo\byte} SRAM, and \SI{128}{\kilo\byte} flash. Three
properties make it a good measurement target:

\begin{itemize}
  \item The Cortex-M3 has a well-documented, deterministic interrupt
        model with a minimum interrupt latency of 12 cycles
        (\SI{167}{\nano\second} at \SI{72}{\mega\hertz})~\cite{armcm3trm}.
  \item The Data Watchpoint and Trace (DWT) unit provides a free-running
        CPU cycle counter (\texttt{DWT\_CYCCNT}), enabling
        cycle-granular self-instrumentation.
  \item The platform is fully supported by both Zephyr and Renode,
        allowing bit-identical firmware to run in both environments.
\end{itemize}

The firmware is built on Zephyr RTOS~3.6.0~\cite{zephyr}, a
widely-deployed open-source RTOS. Zephyr's kernel timing primitives
used in this work are:

\begin{itemize}
  \item \ksleep{} --- suspends the calling thread for a given number of
        ticks. With the default \SI{1000}{\hertz} tick rate, requested
        durations are rounded \emph{up} to the next tick boundary,
        introducing a systematic $0$--$1$\,\ms{} quantisation delay.
  \item \texttt{k\_busy\_wait()} --- spins on the cycle counter without
        yielding, providing sub-tick delays.
  \item \ktimer{} --- hardware-timer-backed periodic callbacks,
        dispatched from the system clock interrupt independent of the
        scheduler tick.
  \item Thread scheduling --- preemptive, priority-based scheduling
        with optional time slicing.
\end{itemize}

\section{Firmware Rehosting}
\label{sec:bg:rehosting}

Rehosting (also \emph{re-hosting} or firmware emulation) denotes the
execution of firmware compiled for an embedded target inside a
software environment on a host machine~\cite{fasano2021sok}. Depending
on how much of the hardware is modelled, approaches range from
high-level API interception (HALucinator~\cite{clements2020halucinator})
over automatically inferred peripheral models
(P$^2$IM~\cite{feng2020p2im}, Pretender~\cite{gustafson2019pretender},
$\mu$Emu~\cite{zhou2021uemu}) to full-system emulation with hand-written
peripheral models (QEMU~\cite{bellard2005qemu}, Renode~\cite{renode}).
Full-system emulation offers the strongest fidelity guarantees for
\emph{functional} behaviour because the unmodified binary executes
against a complete register-level model of the SoC.

\section{The Renode Emulation Framework}
\label{sec:bg:renode}

Renode~\cite{renode} is an open-source emulation framework by Antmicro,
targeted specifically at embedded and multi-node IoT systems. Its
architecture is relevant to this thesis in four respects.

\paragraph{Platform descriptions.} A machine is assembled from a
\texttt{.repl} (Renode platform) file that instantiates peripherals at
bus addresses. Peripherals are C\# classes loaded at runtime; users can
add custom peripherals \emph{without recompiling Renode}. This
extension mechanism is the foundation of the Temporal Fidelity Layer.

\paragraph{Translation-based CPU emulation.} Guest instructions are
executed through a translation library derived from tlibs. Renode does
not simulate microarchitectural effects (pipeline, caches, flash wait
states); instead, each CPU is assigned a nominal performance in MIPS
(here: 72), and executed instruction counts are converted into virtual
time at that fixed rate. Deviation source \srcS{1}
(\cref{sec:problem:sources}) follows directly from this abstraction.

\paragraph{Quantum-based virtual time.} Virtual time advances in
discrete slices called \emph{quanta} (default \SI{100}{\micro\second}).
Each CPU executes up to one quantum's worth of instructions before
synchronising with the emulation's master clock. Timer interrupts,
inter-peripheral events, and cross-machine messages are therefore
aligned to quantum boundaries, which quantises event timing
(source \srcS{4}).

\paragraph{Deterministic execution.} Given identical inputs, a Renode
simulation is deterministic: there is no analog noise, no oscillator
drift (source \srcS{3}), and no physical transmission delay between
emulated nodes (source \srcS{5}). Determinism is desirable for
reproducible debugging, but it is precisely the property that separates
emulated timing from physical timing.

\section{Measurement Instrumentation}
\label{sec:bg:instrumentation}

On physical hardware, firmware-internal events are made externally
observable by toggling GPIO pins, captured with a Saleae Logic~8
analyser~\cite{saleae} at \SI{100}{\mega\sample\per\second} (\SI{10}{\nano\second}
resolution). GPIO instrumentation was chosen over trace units (ITM/ETM)
because the identical mechanism exists in Renode: the emulated GPIO
port is hooked to log state changes with virtual-time timestamps,
yielding structurally identical traces from both worlds. The
instrumentation cost itself is quantified in
\cref{sec:eval:instrumentation}.
